\newpage
\phantomsection
\label{prf:axiomatic-equality-function-congruence}
\LRAProofFor{thm:axiomatic-equality-function-congruence}

\begin{remark*}[Return]
\hyperref[thm:axiomatic-equality-function-congruence]{Return to Theorem}
\end{remark*}

\ProofVaultURL{https://github.com/wsollers/lra-proof-vault/tree/master/volume-i/book-logic/axiomatic-equality/thm-axiomatic-equality-function-congruence}

\begin{theorem*}[Function Congruence]
For any unary function symbol \(f\),
\[
  \forall x\,\forall y\,(x=y\to f(x)=f(y)).
\]
More generally, equal inputs in each coordinate yield equal outputs.
\end{theorem*}

\begin{proof}[Professional Standard Proof]
\LRAProofBodyStart
Assume \(x=y\). By reflexivity of equality, \(f(x)=f(x)\). Let \(\Phi(z)\)
be the formula \(f(x)=f(z)\). Then \(\Phi(x)\) holds. Since \(x=y\),
substitution of identicals applied to the formula-context \(\Phi\) gives
\(\Phi(y)\). Hence \(f(x)=f(y)\). Therefore \(x=y\to f(x)=f(y)\), and
universal generalization gives
\[
  \forall x\,\forall y\,(x=y\to f(x)=f(y)).
\]
\end{proof}

\begin{proof}[Detailed Learning Proof]
\LRAProofBodyStart
\textbf{Step 1.} Assume the hypothesis:
\[
  x=y.
\]
We want to derive \(f(x)=f(y)\).

\textbf{Step 2.} Reflexivity applies to every term, including the term
\(f(x)\). Hence
\[
  f(x)=f(x).
\]

\textbf{Step 3.} Choose the formula-context
\[
  \Phi(z) \quad\text{to mean}\quad f(x)=f(z).
\]
The variable \(z\) marks the input position where substitution is being
performed. The first occurrence \(f(x)\) is held fixed as a parameter of the
chosen formula-context.

\textbf{Step 4.} With this choice of \(\Phi\), the statement \(\Phi(x)\) is
exactly
\[
  f(x)=f(x),
\]
which is true by reflexivity. Since the assumption is \(x=y\), substitution of
identicals transfers the truth of \(\Phi(x)\) to \(\Phi(y)\).

\textbf{Step 5.} But \(\Phi(y)\) is exactly
\[
  f(x)=f(y).
\]
Thus \(f(x)=f(y)\). Since this conclusion was derived from the arbitrary
assumption \(x=y\), we have proved
\[
  x=y\to f(x)=f(y),
\]
and hence
\[
  \forall x\,\forall y\,(x=y\to f(x)=f(y)).
\]
\end{proof}

\begin{remark*}[Proof structure]
The proof uses reflexivity to obtain the true instance \(f(x)=f(x)\), then uses
substitution of identicals along \(x=y\). The formula-context
\(\Phi(z)\equiv f(x)=f(z)\) records that only the second input occurrence is
being replaced.
\end{remark*}

\begin{dependencies}
\begin{itemize}
  \item \hyperref[ax:axiomatic-equality-reflexivity]{Reflexivity of Equality}
  \item \hyperref[ax:axiomatic-equality-substitution]{Substitution of Identicals}
\end{itemize}
\end{dependencies}

\clearpage
